\documentclass[zihao=-4,fontset=ubuntu]{graduateart}
\usepackage{amsmath,amsthm,amssymb,amsfonts}
\definecolor{nuanbai}{HTML}{F5F5F5}
\pagecolor{nuanbai!5}
\usepackage{ninecolors}
\usepackage[colorlinks,
            pdfstartview=FitH,
            CJKbookmarks=true,
            bookmarksnumbered=true,
            bookmarksopen=true,
            linkcolor=black,
            filecolor=purple5,
            anchorcolor=purple5,
            linktocpage=true,
            urlcolor=purple5,
            citecolor=purple5,
            pdfencoding={auto},
			pdfauthor={ Shilong.Lu},pdfsubject={article},
			pdftitle={titlename},pdfkeywords={article},
            breaklinks,
            unicode,
            linktoc=all,
            plainpages=false,
            pdfusetitle,]{hyperref}
\setlength{\headheight}{13.59999pt}
\usepackage{appendix}
\usepackage[sort&compress]{gbt7714}
\citestyle{numbers}% 可选参数含义：存在多篇引文时自动压缩序号

\begin{document}
\author{陆~世~龙}
\teacher{卢~卫~君\quad 教授}
\subject{基~础~数~学}
\academy{数~学~与~物~理~学~院}
\finishedtime{2022年9月}
\grade{2022级}
\subtitle{微分几何中的相关问题研究与说明}
\logo{R-C-removebg-preview.png}
\makecover
\input{undertitlepage}
\input{titlepage}

\pagenumbering{Roman}
\input{abstract}
\input{enabstarct}
\tableofcontents\clearpage

\pagenumbering{arabic}
\section{第一节}
测试一下参考文献\cite{吴大任1979微分几何讲义}
\begin{lstlisting}[language=Python, caption=Python example]
import numpy as np
    
def incmatrix(genl1,genl2):
    m = len(genl1)
    n = len(genl2)
    M = None #to become the incidence matrix
    VT = np.zeros((n*m,1), int)  #dummy variable
    
    #compute the bitwise xor matrix
    M1 = bitxormatrix(genl1)
    M2 = np.triu(bitxormatrix(genl2),1) 

    for i in range(m-1):
        for j in range(i+1, m):
            [r,c] = np.where(M2 == M1[i,j])
            for k in range(len(r)):
                VT[(i)*n + r[k]] = 1;
                VT[(i)*n + c[k]] = 1;
                VT[(j)*n + r[k]] = 1;
                VT[(j)*n + c[k]] = 1;
                
                if M is None:
                    M = np.copy(VT)
                else:
                    M = np.concatenate((M, VT), 1)
                
                VT = np.zeros((n*m,1), int)
    
    return M
\end{lstlisting}
行内\verb|while(1)|

\newpage


\appendix
\section{附录：第一节 测试}






\newpage

\bibliographystyle{gbt7714-numerical}   % 我选择了使用顺序编码制
\bibliography{ref}
\end{document}
